\begin{enunciado}{\ejExtra}
  Hallar el menor $n \en \naturales$ tal que:\par
  \begin{enumerate}[label=\roman*)]
    \item $(n : 2528) = 316$
    \item $n$ tiene exáctamente 48 divisores positivos
    \item $27 \noDivide n$
  \end{enumerate}
\end{enunciado}

Analizo y acomodo todos los números:
$$
  \llave{l}{
    \flecha{factorizo}[2528] 2528 = 2^5 \cdot 79 \Tilde \\
    \flecha{factorizo}[316] 316 = 2^2 \cdot 79 \Tilde
  }
$$
Reescribo las condiciones del enunciado porque pintó:
$$
  (n:2^5 \cdot 79) = 2^2 \cdot 79
$$
Quiero encontrar:
$$
  n = 2^{\alpha_2} \cdot 3^{\alpha_3} \cdot 5^{\alpha_5} \cdot 7^{\alpha_7} \cdots 79^{\alpha_{79}} \cdots
$$
La condición del MCD me obliga a que $n$ tenga ciertos primos con ciertos exponentes:
$$
  (n:2^5 \cdot 79) = 2^2 \cdot 79
  \entonces
  \llave{ll}{
    \alpha_2 = 2,                    & \text{dado que $2^2 \cdot 79 \divideA n$. \ul{Busco el menor $n$}.} \\
    \alpha_{79} \geq 1,              & \text{Todavía no sé exactamente cuánto es $\alpha_{79}$.}
  }
$$
Notar que $\alpha_3 < 3$ si no pasaría que $3^3 = 27 \divideA n$

La estrategia sigue con el primo más chico que haya. Para saber la cantidad de divisores de un número acá está el machete
\hyperlink{teoria-4:cantidadDivisores}{\click}.
$$
  \llave{l}{
    48 = \ub{(\alpha_2 + 1)}{2 + 1} \cdot (\alpha_3 + 1) \cdots                                                                             \\
    48 = 3 \cdot (\alpha_3 + 1) \cdots                                                                                                      \\
    16 = (\alpha_3 + 1) \cdot (\alpha_5 + 1) \cdot (\alpha_7 + 1) \cdots \ub{(\alpha_{79} + 1)}{=2\text{ quiero el menor}} \cdots \\
    8 = (\alpha_3 + 1) \cdot (\alpha_5 + 1) \cdot (\alpha_7 + 1) \cdots                                                                     \\
    8 = \ub{(\alpha_3 + 1)}{=2} \cdot \ub{(\alpha_5 + 1)}{=2} \cdot \ub{(\alpha_7 + 1)}{=2} \cdot 1 \cdots 1                                \\
  }
$$

El $n$ que cumple lo pedido sería $n = 2^2 \cdot 3^1 \cdot 5^1 \cdot 7^1 \cdot 79^1$

\begin{aportes}
  \item \aporte{\dirRepo}{naD GarRaz \github}
\end{aportes}
